\documentclass[11pt, a4paper]{article}

\usepackage[utf8]{inputenc}
\usepackage[T1]{fontenc}
\usepackage{amsmath}
\usepackage{amssymb}
\usepackage{graphicx}
\usepackage{geometry}
\usepackage{fancyhdr}
\usepackage{enumitem}
\usepackage{array}
\usepackage{eso-pic}
\usepackage{esvect}
\usepackage[most]{tcolorbox}
\usepackage{siunitx}
\DeclareSIUnit{\litre}{l}
\sisetup{
    locale = DE,
    inter-unit-product = \cdot,
    per-mode = symbol-or-fraction
}

\makeatletter
\def\input@path{{./}{../../}}
\makeatother

\geometry{a4paper, top=2.5cm, bottom=2.5cm, left=2.5cm, right=2.5cm}

\input{defs}

\definecolor{boxcol_back_lgreen}{RGB}{220, 255, 220}
\definecolor{boxcol_back_white}{RGB}{255, 255, 255}
\definecolor{boxcol_frame_green}{RGB}{30, 180, 30}
\newtcolorbox{solutionbox}[1][] {
  colback=boxcol_back_white,
  colframe=boxcol_frame_green,
  colbacktitle=boxcol_back_lgreen,
  fonttitle=\bfseries,
  coltitle=black,
  boxsep=5pt,
  arc=1mm,
  enhanced,
  attach boxed title to top left={yshift=-2mm, xshift=3mm},
  title=Lösung,
  #1
}

\pagestyle{fancy}
\fancyhf{}
\cfoot{\thepage}

\begin{document}

\section*{\LARGE Lösung: Kurztest 4}
\noindent
\begin{tabular}{ll}
    \textbf{Gruppe: } & \LARGE\textbf{A}  \\[0.3cm]
\end{tabular}

\begin{enumerate}[label=\textbf{\arabic*.},itemsep=0.50cm]

    \item (1 P) - \textbf{Grundgrößen:}
    \begin{solutionbox}
        Mögliche Formeln sind:
        \begin{itemize}[itemsep=0pt, topsep=2pt]
            \item Arbeit: $W = F \cdot s$ mit $[W] = \si{\joule}$
            \item kinetische Energie: $E_{\text{kin}} = \frac{1}{2}mv^2$ mit $[E_{\text{kin}}] = \si{\joule}$
            \item potentielle Energie: $E_{\text{pot}} = mgh$ mit $[E_{\text{pot}}] = \si{\joule}$
            \item Leistung: $P = \frac{W}{t}$ oder $P = Fv$ mit $[P] = \si{\watt}$
        \end{itemize}
    \end{solutionbox}

    \item (1,5 P) - \textbf{Hubarbeit:}
    \begin{solutionbox}
        Es gilt
        \[
        W_{\text{Hub}} = mgh = \SI{15}{\kilogram} \cdot \SI{10}{\meter\per\second\squared} \cdot \SI{1.6}{\meter} = \SI{240}{\joule}
        \]
        \begin{itemize}[itemsep=0pt, topsep=2pt]
            \item Die verrichtete Hubarbeit beträgt $\SI{240}{\joule}$.
            \item Die potentielle Energie nimmt ebenfalls um $\Delta E_{\text{pot}} = \SI{240}{\joule}$ zu.
        \end{itemize}
    \end{solutionbox}

    \item (1,5 P) - \textbf{Freier Fall und Energieerhaltung:}
    \begin{solutionbox}
        \begin{itemize}[itemsep=2pt, topsep=2pt]
            \item Anfangsenergie pro Kilogramm:
            \[
            \frac{E_{\text{pot}}}{m} = gh = \SI{10}{\meter\per\second\squared} \cdot \SI{5}{\meter} = \SI{50}{\joule\per\kilogram}
            \]
            \item Mit Energieerhaltung gilt $mgh = \frac{1}{2}mv^2$. Daraus folgt
            \[
            v = \sqrt{2gh} = \sqrt{2 \cdot 10 \cdot 5} = \SI{10}{\meter\per\second}
            \]
        \end{itemize}
    \end{solutionbox}

    \item (1,5 P) - \textbf{Bremsweg aus der Arbeit:}
    \begin{solutionbox}
        Zuerst Umrechnung der Geschwindigkeit:
        \[
        \SI{72}{\kilo\meter\per\hour} = \frac{72}{3.6} \, \si{\meter\per\second} = \SI{20}{\meter\per\second}
        \]
        Die kinetische Energie zu Beginn ist
        \[
        E_{\text{kin}} = \frac{1}{2}mv_0^2 = \frac{1}{2} \cdot 1200 \cdot 20^2 = \SI{240000}{\joule}
        \]
        Die Bremsarbeit ist $W_B = F_B \cdot s$ und entspricht betragsmäßig der anfänglichen kinetischen Energie:
        \[
        F_B s = E_{\text{kin}}
        \]
        also
        \[
        s = \frac{E_{\text{kin}}}{F_B} = \frac{240000}{4800} = \SI{50}{\meter}
        \]
    \end{solutionbox}

    \item (1,5 P) - \textbf{Mechanische Leistung eines PKW:}
    \begin{solutionbox}
        Die gesamte entgegenwirkende Kraft ist
        \[
        F_{\text{ges}} = F_{LW} + F_R = \SI{320}{\newton} + \SI{180}{\newton} = \SI{500}{\newton}
        \]
        Die Geschwindigkeit beträgt
        \[
        \SI{90}{\kilo\meter\per\hour} = \SI{25}{\meter\per\second}
        \]
        Damit ergibt sich für die Leistung
        \[
        P = F_{\text{ges}} v = \SI{500}{\newton} \cdot \SI{25}{\meter\per\second} = \SI{12500}{\watt} = \SI{12.5}{\kilo\watt}
        \]
    \end{solutionbox}

    \item (1,5 P) - \textbf{Energieübertragung:}
    \begin{solutionbox}
        Aus $E_{\text{kin}} = \frac{1}{2}mv^2$ folgt
        \[
        v = \sqrt{\frac{2E_{\text{kin}}}{m}} = \sqrt{\frac{2 \cdot 3250}{65}} = \sqrt{100} = \SI{10}{\meter\per\second}
        \]
        Am höchsten Punkt ist die gesamte Anfangsenergie in potentielle Energie umgewandelt:
        \[
        E_{\text{kin}} = mgh_{\text{max}}
        \]
        also
        \[
        h_{\text{max}} = \frac{3250}{65 \cdot 10} = \SI{5}{\meter}
        \]
        Die Person erreicht also eine maximale Höhe von $\SI{5}{\meter}$ über dem Absprungpunkt.
    \end{solutionbox}

    \item (1,5 P) - \textbf{Vorzeichen der Arbeit:}
    \begin{solutionbox}
        Arbeit ist positiv, wenn Kraft und Bewegung in dieselbe Richtung zeigen, negativ bei entgegengesetzter Richtung und null bei senkrechter Richtung.
        \begin{itemize}[itemsep=2pt, topsep=2pt]
            \item Positive Arbeit: die Zugkraft
            \item Negative Arbeit: die Reibungskraft
            \item Keine Arbeit: Gewichtskraft und Normalkraft, da beide senkrecht zur horizontalen Bewegung stehen
        \end{itemize}
    \end{solutionbox}

\end{enumerate}

\vspace{1cm}

\end{document}